% !TEX root = ../diplomarbeit.tex

\chapter*{Abstract}
\addcontentsline{toc}{chapter}{Abstract}

\emph{Das Abstract ist eine prägnante Zusammenfassung der gesamten Arbeit in englischer Sprache. Es umfasst typischerweise 150--250 Wörter und beschreibt das Problem, die verwendete Methodik, die wichtigsten Ergebnisse und die Schlussfolgerungen. Das Abstract ermöglicht Lesenden eine schnelle Einschätzung der Relevanz der Arbeit.}

% --------------------------------------------------------
\placeholderbox{Example}{%
Modern software development faces significant challenges in managing dependencies across large-scale projects, with version conflicts consuming substantial developer time. Existing dependency management tools provide automatic updates but fail to consider semantic compatibility at the API level, leading to runtime errors and costly hotfixes.

This thesis presents an intelligent dependency analyzer that leverages static code analysis to detect actual API usage patterns and automatically suggest compatible library versions. The system combines abstract syntax tree analysis with a machine learning classifier trained on 50,000 real-world dependency conflicts.

Evaluation on ten open-source projects (10,000--100,000 lines of code) demonstrates that the analyzer identifies 94\,\% of breaking changes before integration, with an average analysis time of 3.2 minutes per project. The false positive rate was reduced to 8\,\% compared to 23\,\% in existing tools.

The proposed solution significantly reduces integration risks and accelerates development cycles. Future work will extend support to dynamic languages and incorporate runtime behavior analysis.
}

\chapter*{Kurzfassung}
\addcontentsline{toc}{chapter}{Kurzfassung}

\emph{Die Kurzfassung ist eine prägnante Zusammenfassung der gesamten Arbeit in deutscher Sprache. Sie umfasst typischerweise 150--250 Wörter und beschreibt das Problem, die verwendete Methodik, die wichtigsten Ergebnisse und die Schlussfolgerungen. Die Kurzfassung ermöglicht Lesenden eine schnelle Einschätzung der Relevanz der Arbeit.}

\placeholderbox{Beispiel}{%
Die moderne Softwareentwicklung steht vor erheblichen Herausforderungen bei der Verwaltung von Abhängigkeiten in großangelegten Projekten, wobei Versionskonflikte erhebliche Entwicklerzeit binden. Bestehende Dependency-Management-Tools bieten zwar automatische Updates, berücksichtigen jedoch nicht die semantische Kompatibilität auf API-Ebene, was zu Laufzeitfehlern und kostspieligen Hotfixes führt.

Diese Diplomarbeit präsentiert einen intelligenten Dependency-Analyzer, der mittels statischer Code-Analyse tatsächliche API-Nutzungsmuster erfasst und automatisch kompatible Bibliotheksversionen vorschlägt. Das System kombiniert Abstract-Syntax-Tree-Analyse mit einem Machine-Learning-Klassifikator, der mit 50.000 realen Abhängigkeitskonflikten trainiert wurde.

Die Evaluierung an zehn Open-Source-Projekten (10.000--100.000 Zeilen Code) zeigt, dass der Analyzer 94\,\% der Breaking Changes vor der Integration identifiziert, bei einer durchschnittlichen Analysezeit von 3,2 Minuten pro Projekt. Die False-Positive-Rate wurde gegenüber bestehenden Tools von 23\,\% auf 8\,\% reduziert.

Die vorgeschlagene Lösung reduziert Integrationsrisiken signifikant und beschleunigt Entwicklungszyklen. Zukünftige Arbeiten werden die Unterstützung dynamischer Sprachen erweitern und Laufzeitverhaltensanalyse integrieren.
}